\clearpage
\section{From text to tokens}\label{sec:tokens}
\subsection*{Text processing makes modeling decisions}
A tokenizer maps text to integer IDs; its decoder maps those IDs back to
text. The complete pipeline includes earlier decisions about decoding,
normalization, and splitting. A later decoder cannot reconstruct distinctions
that an earlier step has discarded. Split documents into training,
development, and test sets before fitting vocabulary or preprocessing choices.

\fig{text-pipeline}{Every arrow has a contract. Normalization may intentionally
lose information; lossless tokenization applies to the text after that step.}

A visible symbol, a Unicode code point, and a UTF-8 byte are different
units. The precomposed letter \texttt{\'e} is one code point and two bytes.
The visually similar sequence \texttt{e} followed by a combining acute accent
uses two code points and three bytes. Rendering alone cannot identify the
stored sequence.

\fig{unicode-units}{Two encodings of an accented letter. Hexadecimal byte values
describe storage; grapheme appearance describes what a reader sees.}

NFC normalization can compose canonically equivalent sequences. NFKC also
applies compatibility mappings, which may remove distinctions useful to a
task. Lowercasing is another transformation: it maps both \texttt{A} and
\texttt{a} to the same string. None is a universal requirement for modeling.
Document the intended transformation and evaluate on text processed the same way.

\fig{boundaries}{Extraction can omit punctuation and spaces. A lossless partition
covers the complete string and can also constrain which adjacent pieces may merge.}

For the input \texttt{hi, 3!}, extracting alphabetic runs produces only
\texttt{hi}. A partition into \texttt{hi}, \texttt{,}, a space, \texttt{3},
and \texttt{!} preserves every character. A regular expression can implement
either behavior; the word ``regex'' alone says nothing about reversibility.

\selfcheck{If a tokenizer round-trips lowercased text, does it also recover the
original capitalization? No: test the complete preprocessing pipeline separately.}

\nextpage{Learning a vocabulary with byte-pair encoding}
Word occurrences count positions, while word types count distinct forms.
For \texttt{low low lower}, whitespace splitting gives three occurrences
and two types; UTF-8 storage uses 13 bytes, including two spaces. Token IDs
are arbitrary labels, not numerical measurements of meaning.

Byte-pair encoding (BPE) begins from a base alphabet and repeatedly adds a
token formed by merging a frequent adjacent pair. In the course's byte-level
example the base vocabulary contains all 256 byte values. Each training
sequence has a frequency weight; boundaries prevent merges between sequences.
The toy corpus is \texttt{low}$\times5$ and \texttt{lower}$\times2$.

\fig{bpe-merges}{Two weighted merges. Both initial winning pairs have count seven;
the course's deterministic tie rule selects \texttt{l o} first.}

The initial weighted length is $5(3)+2(5)=25$. The pair
$(\texttt{l},\texttt{o})$ occurs once in each of seven weighted words, so
its merge reduces the total to 18. Merging $(\texttt{lo},\texttt{w})$
reduces it to 11. The new IDs are 256 for \texttt{lo} and 257 for
\texttt{low}; two merges give 258 vocabulary entries.

\textbf{Training algorithm.}
\begin{enumerate}
\item Count adjacent pairs inside each current sequence and multiply by
the sequence's corpus frequency.
\item Select the most frequent pair, resolving ties deterministically.
Assign its concatenation a new token ID and a merge rank.
\item Replace non-overlapping occurrences of that pair, then recount.
Stop after the requested number of merges or when no eligible pair remains.
\end{enumerate}
Counting and replacement are different operations. In \texttt{aaa}, the
pair \texttt{a a} occurs at two overlapping positions, but one replacement
can consume only two of the three symbols. For \texttt{aaa}$\times2$ and
\texttt{bc}$\times2$, its weighted pair count is four while the token total
decreases only from ten to eight.

\takeaway{Vocabulary growth is learned from training text. Deterministic
ties, allowed boundaries, and overlap handling are part of the algorithm.}
\selfcheck{Recompute the second merge counts after the first replacement.
Why would merging across spaces change what the vocabulary can represent?}
\textbf{Reading connection.} \readingref{R03} applies BPE to subword units
within words. Its character-based setup differs from this byte-level example.

\nextpage{Encoding new text and choosing a vocabulary}
Encoding freezes the learned merge list. Begin with the new text's base
symbols, find an eligible adjacent pair with the earliest learned rank,
merge it, and repeat. Do not count new-text frequencies to learn new rules:
the model's embedding rows already correspond to a fixed vocabulary.

\fig{frozen-bpe}{The learned \texttt{lo} and \texttt{low} merges apply to
\texttt{lowest}; the remaining bytes have no learned merges in this example.}

The resulting IDs are $[257,101,115,116]$. Decoding concatenates each
token's byte string and then decodes the complete UTF-8 byte stream.
A single token can end inside a multibyte character, so decoding tokens
independently may fail. Byte coverage ensures that valid input bytes can be
represented; it does not ensure efficient segmentation for every language.
Special tokens such as \BOS{} and \EOS{} have separate control meanings
and need an explicit policy when similar strings appear in ordinary text.

\fig{vocabulary-tradeoff}{Left: larger vocabularies increase table rows while often
shortening sequences. Right: held-out Chinese token counts from the L01 toy
experiment, using the same two sentences at every merge budget.}

At 0, 8, and 32 merges, mixed English/Chinese training gives 54, 48, and 35
tokens on those Chinese sentences. English-only training gives 54 each time.
The tokenizer still covers the text, but its merges offer no compression there.
These small constructed corpora demonstrate a mechanism; they do not rank
tokenizers on broad language-model quality.

For a table with $K$ rows and width $d$, storage grows as $Kd$. A materialized
output distribution at $BT$ positions has $BTK$ entries. Shorter sequences
may reduce position-dependent work, while a larger vocabulary increases
table and readout costs. The resulting tradeoff depends on the model and
text distribution; token compression alone cannot settle it.

\takeaway{Fit the tokenizer on training data, freeze it for evaluation, and
report both representation behavior and downstream modeling results.}
\selfcheck{Why can a tokenizer with fewer tokens have higher per-token perplexity?
Section 4 will separate the prediction unit from the amount of text scored.}
